	\chapter{Le cadenze e la sintassi musicale}
	\label{chap:cadenze}
	
	Se gli accordi sono le parole del linguaggio musicale, le \textbf{cadenze} ne sono la punteggiatura. 
		
\begin{definizione}[La Cadenza]
	\label{def:cadenza}
	Una \textbf{cadenza} è una formula melodico-armonica costituita da una successione di due o più accordi che crea un punto di articolazione, pausa o conclusione nel discorso musicale.
\end{definizione}

Le cadenze svolgono una funzione analoga alla punteggiatura nel linguaggio verbale: delimitano le frasi, creano momenti di sospensione o di chiusura, e strutturano il flusso temporale della composizione.

Le cadenze possono essere classificate in base al loro effetto percettivo e alla loro funzione formale. Alcune cadenze producono un senso di \textbf{conclusione} e \textbf{riposo}, altre creano \textbf{sospensione} e \textbf{attesa}. Questa capacità di guidare le aspettative dell'ascoltatore è uno degli strumenti più potenti a disposizione del compositore per dare forma e significato alla musica.

Storicamente, le cadenze sono state codificate nella pratica compositiva come formule ricorrenti, riconoscibili e immediatamente efficaci. La loro forza non risiede solo nella successione di accordi in sé, ma anche nella condotta delle voci, nel ritmo armonico e nella posizione metrica in cui vengono collocate.
	
	\section{Cadenze conclusive}
	\label{sec:cadenze_conclusive}
	
	Le cadenze conclusive, come suggerisce il nome, producono un effetto di chiusura e riposo. Sono utilizzate per terminare le frasi musicali in modo definitivo.
	
	\begin{definizione}[Cadenza Perfetta]
		\label{def:cadenza_perfetta}
		La \textbf{cadenza perfetta} (o autentica perfetta) è la formula conclusiva più forte e inequivocabile del sistema tonale. Consiste nella successione degli accordi \textbf{V - I}, entrambi in stato fondamentale. Per essere considerata "perfetta", deve soddisfare una condizione aggiuntiva: la nota finale al soprano deve essere la tonica.
	\end{definizione}
	
	La cadenza perfetta è il culmine della tensione tonale: la dominante, con la sua sensibile che spinge verso l'alto, risolve nel centro di gravità della tonica, portando il discorso a una quiete completa.
	
	\begin{esempio}[Cadenza Perfetta in Do Maggiore]
		\label{ex:cadenza_perfecta}
		\musica{
			\partilate
			{ \key do \major \time 4/4 }
			{ re'2 do'2 }
			{ si2 sol2 }
			{ sol2 mi2 }
			{ sol,2 do2 }
			\lyricmode { "V"2 "I"2 }
			\figuremode {}
		}
	\end{esempio}

	\begin{definizione}[Cadenza Imperfetta]
		\label{def:cadenza_imperfetta}
		La \textbf{cadenza imperfetta} (o autentica imperfetta) è una cadenza V-I (o V7-I) in cui almeno una delle seguenti condizioni non è soddisfatta:
		\begin{itemize}
			\item L'accordo di V non è in stato fondamentale (es. V6 o V6/5 invece di V).
			\item L'accordo di I non è in stato fondamentale (es. I6 invece di I).
			\item Il soprano all'accordo di arrivo non ha la tonica (es. ha la terza o la quinta).
		\end{itemize}
		Produce un senso di chiusura meno definitivo della cadenza perfetta, ed è spesso usata a metà periodo (\ref{chap:struttura_frase}).
	\end{definizione}

	\begin{esempio}[Esempi di Cadenza Imperfetta]
		Mostriamo due varianti: soprano sulla terza e V6 invece di V.
		\musica{
			\partilate
			{ \key do \major \time 4/4 }
			{ re'2 mi' | r2 re' | si2 do' }
			{ si2 do' | r2 si | sol2 sol }
			{ sol2 sol | r2 fa | re2 mi }
			{ sol,2 do | r2 sol, | sol,2 do }
			\lyricmode { "V-I (sop.3a)" | "V6-I" }
			\figuremode { s2 s | r <6> | s s }
		}
	\end{esempio}

    \begin{regola}[Gerarchia delle cadenze]
        \label{reg:gerarchia_cadenze}
        La forza conclusiva delle cadenze segue questa gerarchia decrescente:
        \begin{center}
            \textbf{C. Perfetta > C. Imperfetta > C. Plagale > C. Sospesa > C. d'Inganno}
        \end{center}
        Questa gerarchia guida il compositore nella scelta della cadenza in base alla collocazione nella struttura formale (\ref{chap:struttura_frase}).
    \end{regola}
	
	\begin{definizione}[Cadenza Plagale]
		\label{def:cadenza_plagale}
		La \textbf{cadenza plagale} consiste nella successione \textbf{IV - I}. Produce un senso di chiusura più debole e statico rispetto alla cadenza perfetta. È spesso usata come formula conclusiva aggiuntiva dopo una cadenza perfetta, ed è universalmente nota come la cadenza dell'"Amen" nella musica liturgica.
	\end{definizione}
	
	\begin{esempio}[Cadenza Plagale in Do Maggiore]
		\label{ex:cadenza_plagale}
		\musica{
			\partilate
			{ \key do \major \time 4/4 }
			{ fa'2 mi'2 }
			{ do'2 do'2 }
			{ la2 sol2 }
			{ fa,2 do2 }
			\lyricmode { "IV"2 "I"2 }
			\figuremode {}
		}
	\end{esempio}
	
	\section{Cadenze sospensive}
	\label{sec:cadenze_sospensive}
	
	Le cadenze sospensive creano un effetto di interruzione e attesa, lasciando il discorso musicale "aperto". Sono l'equivalente musicale di una virgola o di un punto interrogativo.
	
	\begin{definizione}[Cadenza Sospesa]
		\label{def:cadenza_sospesa}
		Una \textbf{cadenza sospesa} (o mezza cadenza) è una qualsiasi formula armonica che termina sull'accordo di \textbf{Dominante (V)}. L'accordo che la precede può essere la Tonica (I), la Sottodominante (IV), la Sopratonica (II) o altri. La sua funzione è quella di creare un'aspettativa per la risoluzione sulla tonica, che però viene posticipata.
	\end{definizione}
	
	\begin{esempio}[Cadenza Sospesa (IV - V)]
		\label{ex:cadenza_sospesa}
		\musica{
			\partilate
			{ \key do \major \time 4/4 }
			{ fa'2 sol'2 }
			{ do'2 re'2 }
			{ la2 si2 }
			{ fa,2 sol,2 }
			\lyricmode { "IV"2 "V"2 }
			\figuremode {}
		}
	\end{esempio}
	
	\begin{definizione}[Cadenza d'Inganno]
		\label{def:cadenza_inganno}
		La \textbf{cadenza d'inganno} è una delle figure più espressive dell'armonia. Consiste nella successione \textbf{V - VI}. Essa "inganna" l'ascoltatore, che si aspetterebbe la risoluzione sulla tonica (I) dopo la dominante. Al suo posto, appare l'accordo del VI grado, che condivide due note con la tonica ma ne sposta il centro tonale, creando una sorpresa e prolungando la frase musicale.
	\end{definizione}
	
	Nella cadenza d'inganno, la condotta delle voci segue regole precise: la sensibile (terza della dominante) sale comunque alla tonica, mentre il basso sale di grado congiunto. Per evitare quinte parallele, si raddoppia tipicamente la terza dell'accordo del VI grado.
	
	\begin{esempio}[Cadenza d'Inganno in Do Maggiore]
		\label{ex:cadenza_inganno}
		\musica{
			\partilate
			{ \key do \major \time 4/4 }
			{ si'2 do''2 } % La sensibile (Si) sale alla tonica (Do)
			{ re'2 mi'2 }
			{ sol2 la2 }
			{ sol,2 la,2 }   % Il basso sale di grado
			\lyricmode { "V"2 "VI"2 }
			\figuremode {}
		}
	\end{esempio}
	
	\section{La fraseologia}
	\label{sec:fraseologia}
	
	Le cadenze non sono eventi isolati, ma i punti di arrivo di un percorso logico governato dalle \textbf{funzioni armoniche}. La struttura più comune di una frase musicale può essere riassunta dal modello T-S-D-T.
	
	\begin{regola}[Logica Funzionale: T-S-D-T]
		\label{reg:tsdt}
		Il modello \textbf{T-S-D-T} descrive il ciclo di tensione e risoluzione che sta alla base della sintassi tonale:
		\begin{itemize}
			\item \textbf{T (Tonica):} Area di stabilità e riposo (I grado).
			\item \textbf{S (Sottodominante):} Area di allontanamento dalla tonica (IV, II grado).
			\item \textbf{D (Dominante):} Area di massima tensione, che prepara il ritorno alla tonica (V grado).
			\item \textbf{T (Tonica):} Ritorno al punto di riposo, chiusura del ciclo.
		\end{itemize}
		Questa progressione fondamentale (Stabilità $\rightarrow$ Allontanamento $\rightarrow$ Tensione $\rightarrow$ Risoluzione) costituisce l'ossatura di innumerevoli brani musicali.
	\end{regola}
	
	\begin{esempio}[Una frase completa basata sul modello T-S-D-T]
		\label{ex:tsdt}
		Una semplice realizzazione del modello I - IV - V - I.
		\musica{
			\partilate
			{ \key do \major \time 4/4 }
			{ mi'4 fa'4 sol'4 do'4 }
			{ do'4 do'4 re'4 mi'4 }
			{ sol4 la4 si4 sol4 }
			{ do4 fa,4 sol,4 do4 }
			\lyricmode { "I (T)"4 "IV (S)"4 "V (D)"4 "I (T)"4 }
			\figuremode {}
		}
	\end{esempio}
	
	
	% ============================================================
	% CAPITOLO 6
	% ============================================================
